\chapter{1977年北京卷（理）}

\section{解答题}

\begin{problem}
解方程：\(\sqrt{x-1}=3-x\).
\end{problem}

\begin{solution}
由根式有意义且等式右边非负，得 \(x-1\geqslant0\) 且 \(3-x\geqslant0\)，所以 \(1\leqslant x\leqslant3\). 两边平方，得 \(x-1=\left(3-x\right)^2\)，即 \(x^2-7x+10=0\)，所以 \(\left(x-2\right)\left(x-5\right)=0\)，得 \(x=2\) 或 \(x=5\). 其中 \(x=5\) 不在 \(\left[1,3\right]\) 内，是增根；将 \(x=2\) 代回原方程得 \(\sqrt{1}=1\)，成立. 故原方程的解为 \(x=2\).
\end{solution}

\begin{problem}
计算：\(2^{-\frac{1}{2}}+\frac{2^0}{\sqrt{2}}+\frac{1}{\sqrt{2}-1}\).
\end{problem}

\begin{solution}
原式 \(=\dfrac{1}{\sqrt{2}}+\dfrac{1}{\sqrt{2}}+\dfrac{1}{\sqrt{2}-1}=\sqrt{2}+\dfrac{\sqrt{2}+1}{\left(\sqrt{2}-1\right)\left(\sqrt{2}+1\right)}=\sqrt{2}+\sqrt{2}+1=2\sqrt{2}+1\).
\end{solution}

\begin{problem}
已知 \(\lg 2=0.3010\)，\(\lg 3=0.4771\)，求 \(\lg\sqrt{45}\).
\end{problem}

\begin{solution}
因为 \(45=5\times3^2\)，且 \(\lg5=\lg10-\lg2=1-0.3010=0.6990\)，所以 \(\lg45=\lg5+2\lg3=0.6990+2\times0.4771=1.6532\). 因此 \(\lg\sqrt{45}=\dfrac{1}{2}\lg45=\dfrac{1}{2}\times1.6532=0.8266\).
\end{solution}

\begin{problem}
证明：\(\left(1+\tan\alpha\right)^2=\dfrac{1+\sin2\alpha}{\cos^2\alpha}\).
\end{problem}

\begin{solution}
在等式有意义的条件 \(\cos\alpha\neq0\) 下，\(1+\sin2\alpha=1+2\sin\alpha\cos\alpha=\sin^2\alpha+2\sin\alpha\cos\alpha+\cos^2\alpha=\left(\sin\alpha+\cos\alpha\right)^2\). 所以 \(\dfrac{1+\sin2\alpha}{\cos^2\alpha}=\left(\dfrac{\sin\alpha+\cos\alpha}{\cos\alpha}\right)^2=\left(1+\tan\alpha\right)^2\)，原等式成立.
\end{solution}

\begin{problem}
求过两直线 \(x+y-7=0\) 和 \(3x-y-1=0\) 的交点且过 \(\left(1,1\right)\) 点的直线方程.
\end{problem}

\begin{solution}
两直线交点 \(\left(x,y\right)\) 满足
\[
\begin{cases}
x+y-7=0,\\
3x-y-1=0.
\end{cases}
\]
两式相加得 \(4x-8=0\)，所以 \(x=2\)，代入第一式得 \(y=5\). 所求直线过 \(\left(2,5\right)\) 和 \(\left(1,1\right)\)，其斜率为 \(\dfrac{5-1}{2-1}=4\)，故 \(y-1=4\left(x-1\right)\)，即所求直线方程为 \(y=4x-3\).
\end{solution}

\begin{problem}
某工厂今年七月份的产值为 \(100\text{ 万元}\)，以后每月产值比上月增加 \(20\%\)，问今年七月份到十月份总产值是多少？
\end{problem}

\begin{solution}
七月份至十月份的产值依次为 \(100\text{ 万元}\)、\(100\times1.2=120\text{ 万元}\)、\(100\times1.2^2=144\text{ 万元}\) 和 \(100\times1.2^3=172.8\text{ 万元}\). 因而总产值为 \(100+120+144+172.8=536.8\text{ 万元}\).
\end{solution}

\begin{problem}
已知二次函数 \(y=x^2-6x+5\).

\begin{enumerate}
    \item 求出它的图象的顶点坐标和对称轴方程；
    \item 画出它的图象；
    \item 分别求出它的图象和 \(x\) 轴，\(y\) 轴的交点坐标.
\end{enumerate}
\end{problem}

\begin{solution}
\begin{enumerate}
    \item 配方得 \(y=x^2-6x+5=\left(x-3\right)^2-4\)，所以抛物线开口向上，顶点坐标为 \(\left(3,-4\right)\)，对称轴方程为 \(x=3\).
    \item 由顶点 \(\left(3,-4\right)\)、对称轴 \(x=3\) 以及下列交点描点，作出开口向上的抛物线 \(y=\left(x-3\right)^2-4\).
    \item 与 \(x\) 轴的交点满足 \(x^2-6x+5=0\)，即 \(\left(x-1\right)\left(x-5\right)=0\)，所以交点为 \(\left(1,0\right)\) 和 \(\left(5,0\right)\). 令 \(x=0\)，得 \(y=5\)，所以与 \(y\) 轴的交点为 \(\left(0,5\right)\).
\end{enumerate}
\end{solution}

\begin{problem}
一只船以 \(20\text{ 海里/小时}\) 的速度向正东航行，起初船在 \(A\) 处看见一灯塔 \(B\) 在船的北 \(45^\circ\) 东方向，一小时后船在 \(C\) 处看见这个灯塔在船的北 \(15^\circ\) 东方向，求这时船和灯塔的距离 \(CB\).
\end{problem}

\begin{solution}
船向正东航行一小时，所以 \(AC=20\text{ 海里}\). 在三角形 \(ABC\) 中，\(\angle BAC=45^\circ\). 从 \(C\) 看，射线 \(CB\) 为北偏东 \(15^\circ\)，而射线 \(CA\) 指向正西，因此 \(\angle ACB=180^\circ-\left(90^\circ-15^\circ\right)=105^\circ\)，从而 \(\angle ABC=180^\circ-45^\circ-105^\circ=30^\circ\). 由正弦定理，\(\dfrac{CB}{\sin45^\circ}=\dfrac{AC}{\sin30^\circ}\)，所以 \(CB=20\times\dfrac{\sin45^\circ}{\sin30^\circ}=20\sqrt{2}\text{ 海里}\)，约为 \(28.28\text{ 海里}\).
\end{solution}

\begin{problem}
有一个圆内接三角形 \(ABC\)，\(\angle A\) 的平分线交 \(BC\) 于 \(D\)，交外接圆于 \(E\)，求证：\(AD\cdot AE=AC\cdot AB\).
\end{problem}

\begin{solution}
因为 \(AD\) 是 \(\angle A\) 的平分线，由角平分线定理得 \(\dfrac{BD}{DC}=\dfrac{AB}{AC}\). 在三角形 \(ABC\) 中应用斯特瓦尔特定理，得 \(AC^2\cdot BD+AB^2\cdot DC=BC\left(AD^2+BD\cdot DC\right)\). 由 \(\dfrac{BD}{DC}=\dfrac{AB}{AC}\)，可设 \(BD=kAB\)，\(DC=kAC\)，于是 \(BC=k\left(AB+AC\right)\)，且 \(AC^2\cdot BD+AB^2\cdot DC=kAB\cdot AC\left(AB+AC\right)=BC\cdot AB\cdot AC\). 因此 \(AD^2+BD\cdot DC=AB\cdot AC\).

点 \(D\) 在弦 \(BC\) 的内部，所以由相交弦定理，\(DA\cdot DE=DB\cdot DC\). 又因 \(A,D,E\) 三点依次在同一直线上，\(AE=AD+DE\)，从而 \(AD\cdot AE=AD\left(AD+DE\right)=AD^2+AD\cdot DE=AD^2+BD\cdot DC=AB\cdot AC\). 即证得 \(AD\cdot AE=AC\cdot AB\).
\end{solution}

\begin{problem}
当 \(m\) 取哪些值时，直线 \(y=x+m\) 与椭圆 \(\dfrac{x^2}{16}+\dfrac{y^2}{9}=1\) 有一个交点？ 有两个交点？ 没有交点？ 当它们有一个交点时，画出它的图象.
\end{problem}

\begin{solution}
把 \(y=x+m\) 代入椭圆方程，得 \(\dfrac{x^2}{16}+\dfrac{\left(x+m\right)^2}{9}=1\)，即 \(25x^2+32mx+16m^2-144=0\). 其判别式为 \(\Delta=\left(32m\right)^2-4\times25\left(16m^2-144\right)=576\left(25-m^2\right)\).

当 \(\Delta=0\)，即 \(m=\pm5\) 时，直线与椭圆相切，有一个交点；当 \(\Delta>0\)，即 \(-5<m<5\) 时，有两个交点；当 \(\Delta<0\)，即 \(m<-5\) 或 \(m>5\) 时，没有交点.

当 \(m=5\) 时，重根为 \(x=-\dfrac{32m}{2\times25}=-\dfrac{16}{5}\)，切点为 \(\left(-\dfrac{16}{5},\dfrac{9}{5}\right)\)，相应的切线为 \(y=x+5\). 当 \(m=-5\) 时，切点为 \(\left(\dfrac{16}{5},-\dfrac{9}{5}\right)\)，相应的切线为 \(y=x-5\). 作图时分别作出这两条直线，并标出它们与椭圆的上述切点.
\end{solution}

\section{附加题}

\begin{problem}
\begin{enumerate}
    \item 求函数 \(f(x)=\begin{cases}
    x^2\sin\dfrac{\pi}{x}, & x\ne0,\\
    0, & x=0
    \end{cases}\) 的导数.
    \item 求椭圆 \(\dfrac{x^2}{a^2}+\dfrac{y^2}{b^2}=1\) 绕 \(x\) 轴旋转而成的旋转体体积.
\end{enumerate}
\end{problem}

\begin{solution}
\begin{enumerate}
    \item 当 \(x\ne0\) 时，由乘积法则和复合函数求导法则，\(f'(x)=2x\sin\dfrac{\pi}{x}+x^2\cos\dfrac{\pi}{x}\left(-\dfrac{\pi}{x^2}\right)=2x\sin\dfrac{\pi}{x}-\pi\cos\dfrac{\pi}{x}\). 当 \(x=0\) 时，按导数定义，\(f'(0)=\displaystyle\lim_{h\to0}\dfrac{f(h)-f(0)}{h}=\lim_{h\to0}h\sin\dfrac{\pi}{h}=0\)，因为 \(\left|h\sin\dfrac{\pi}{h}\right|\leqslant|h|\to0\). 因此 \(f'(x)=\begin{cases}
    2x\sin\dfrac{\pi}{x}-\pi\cos\dfrac{\pi}{x}, & x\ne0,\\
    0, & x=0.
    \end{cases}\)
    \item 设 \(a>0\)，\(b>0\). 由椭圆方程得，在 \(-a\leqslant x\leqslant a\) 时，\(y^2=b^2\left(1-\dfrac{x^2}{a^2}\right)\). 绕 \(x\) 轴旋转后，横坐标为 \(x\) 处的截面是半径为 \(y\) 的圆，面积为 \(\pi y^2=\pi b^2\left(1-\dfrac{x^2}{a^2}\right)\). 所以旋转体体积为 \(V=\pi\displaystyle\int_{-a}^{a}b^2\left(1-\dfrac{x^2}{a^2}\right)\mathrm{d}x=\pi b^2\left[x-\dfrac{x^3}{3a^2}\right]_{-a}^{a}=\dfrac{4}{3}\pi ab^2\).
\end{enumerate}
\end{solution}

\begin{problem}
\begin{enumerate}
    \item 试用 \(\varepsilon-\delta\) 语言叙述“函数 \(f(x)\) 在点 \(x=x_0\) 处连续”的定义.
    \item 试证明：若 \(f(x)\) 在点 \(x=x_0\) 处连续，且 \(f(x_0)>0\)，则存在一个 \(x_0\) 的邻域 \(\left(x_0-\delta,x_0+\delta\right)\)，在这个邻域内，处处有 \(f(x)>0\).
\end{enumerate}
\end{problem}

\begin{solution}
\begin{enumerate}
    \item 函数 \(f(x)\) 在点 \(x=x_0\) 处连续，是指对任意 \(\varepsilon>0\)，都存在 \(\delta>0\)，使得当 \(\left|x-x_0\right|<\delta\) 时，有 \(\left|f(x)-f(x_0)\right|<\varepsilon\).
    \item 取 \(\varepsilon=\dfrac{f(x_0)}{2}>0\). 由连续性的定义，存在 \(\delta>0\)，使得当 \(\left|x-x_0\right|<\delta\) 时，\(\left|f(x)-f(x_0)\right|<\dfrac{f(x_0)}{2}\). 因而 \(f(x)>f(x_0)-\dfrac{f(x_0)}{2}=\dfrac{f(x_0)}{2}>0\). 所以在邻域 \(\left(x_0-\delta,x_0+\delta\right)\) 内处处有 \(f(x)>0\).
\end{enumerate}
\end{solution}
